% GENERATED FILE. Edit canonical lessons, then run npm run generate:books.
% canonical-source-hash: fnv1a64:77977d90461b6ffe
% canonical-lessons: ES-C305-nombre, ES-C305-apellido, ES-C305-senor, ES-C305-senora, ES-C305-encantado

\chapter{Naming, Properly --- Name, Surname, Sir, Madam, Delighted}
\label{ch:nombre-y-apellido}
\hlchaptermodality{305}

\begin{chapteropening}
\textbf{By the end of this chapter:} \emph{I can ask for and give a first name and a surname in Spanish, address a stranger as senor or senora, and answer an introduction.}

Complete ES-C305-encantado: give both names, address the other person correctly, and answer the introduction.
\end{chapteropening}

\section[el nombre]{el nombre — the word that is also English}

\label{lesson:ES-C305-nombre}



\begin{quote}
Two from last time, before anything new is added.
\end{quote}

\begin{sounds}
\begin{itemize}
\raggedright
  \item Say it once slowly, then at speed. $\to$ pronunciation reference
\end{itemize}
\end{sounds}

\begin{cousinweb}
\textbf{el nombre} — \textquotedblleft{}the name.\textquotedblright{}

Latin \textbf{nomen}. English \emph{name}. These are not borrowings in either direction — they are \textbf{cousins}, both descended from the same ancient word, which is why they still agree in sense after several thousand years apart.

Latin lent the borrowed forms too: a \emph{noun} is a naming word, to \emph{nominate} is to name somebody for a post, and \emph{renown} is being re-named, over and over, by people who were never there.

You already say \emph{me llamo} — \textquotedblleft{}I call myself.\textquotedblright{} \emph{El nombre} is the thing you call yourself \textbf{with}.
\end{cousinweb}

\begin{grammarlens}[title={What you've built}]
A word for the thing your first sentence was already handing over.
\end{grammarlens}

\subsection*{Your turn}
Say these aloud:

\begin{itemize}
\raggedright
  \item \emph{el nombre}
  \item \emph{el nombre, por favor}
  \item \emph{me llamo… ¿y el nombre de usted?}
\end{itemize}

\begin{tcolorbox}[breakable,colback=teal!4,colframe=teal!35!black,title={Before you move on}]
Are \emph{nombre} and \emph{name} borrowed from each other? (No — cousins, from the same ancient word.) What is \emph{renown}? (Being named again and again.)
\end{tcolorbox}
\section[el apellido]{el apellido — the name you are called by}

\label{lesson:ES-C305-apellido}



\begin{quote}
Two from last time, before anything new is added.
\end{quote}

\begin{sounds}
\begin{itemize}
\raggedright
  \item Say it once slowly, then at speed. $\to$ pronunciation reference
\end{itemize}
\end{sounds}

\begin{cousinweb}
\textbf{el apellido} — \textquotedblleft{}the surname,\textquotedblright{} the family name.

It comes from Latin \textbf{appellitare}, the repeating form of \emph{appellare}, \textquotedblleft{}to call on.\textquotedblright{} Not the name you are given — the name you are \textbf{called by}, again and again, by everybody.

English has the same verb in \emph{appeal} (to call on someone) and \emph{appellation} (what a thing is called).

Worth knowing before anyone asks for it: Spanish speakers usually carry \textbf{two} \emph{apellidos}, the father's then the mother's. A form with one box for \textquotedblleft{}surname\textquotedblright{} is a form built for a different country.
\end{cousinweb}

\begin{grammarlens}[title={What you've built}]
Two halves of a full name now, and a warning about forms.
\end{grammarlens}

\subsection*{Your turn}
Say these aloud:

\begin{itemize}
\raggedright
  \item \emph{el apellido}
  \item \emph{el nombre y el apellido}
  \item \emph{el apellido, por favor}
\end{itemize}

\begin{tcolorbox}[breakable,colback=teal!4,colframe=teal!35!black,title={Before you move on}]
What does \emph{appellitare} mean? (To call on again and again.) How many \emph{apellidos} does a Spanish speaker usually have? (Two.)
\end{tcolorbox}
\section[señor]{señor — respect made out of age}

\label{lesson:ES-C305-senor}



\begin{quote}
Two from last time, before anything new is added.
\end{quote}

\begin{sounds}
\begin{itemize}
\raggedright
  \item Say it once slowly, then at speed. $\to$ pronunciation reference
\end{itemize}
\end{sounds}

\begin{cousinweb}
\textbf{señor} — \textquotedblleft{}sir,\textquotedblright{} \textquotedblleft{}Mr.\textquotedblright{}

Latin \textbf{senior} meant \textquotedblleft{}older,\textquotedblright{} the comparative of \emph{senex}, \textquotedblleft{}old.\textquotedblright{} English took it raw as \emph{senior}, and took the French descendant \emph{sieur} — which is inside \emph{monsieur}, \textquotedblleft{}my lord,\textquotedblright{} and worn all the way down to English \emph{sir}.

So \emph{señor}, \emph{sir}, \emph{senior} and \emph{monsieur} are one word, spread out.

The \textbf{ñ} is the letter you met in the writing chapter. And this is what you say to get a stranger's attention politely, ahead of any of last chapter's repairs.
\end{cousinweb}

\begin{grammarlens}[title={What you've built}]
A way to address the person before you ask them anything.
\end{grammarlens}

\subsection*{Your turn}
Say these aloud:

\begin{itemize}
\raggedright
  \item \emph{señor}
  \item \emph{señor, por favor}
  \item \emph{señor — ¿qué significa?}
\end{itemize}

\begin{tcolorbox}[breakable,colback=teal!4,colframe=teal!35!black,title={Before you move on}]
What did \emph{senior} mean? (Older.) Which English one-syllable word is the same root, worn down? (\emph{Sir}.)
\end{tcolorbox}
\section[señora]{señora — the same word, turned}

\label{lesson:ES-C305-senora}



\begin{quote}
Two from last time, before anything new is added.
\end{quote}

\begin{sounds}
\begin{itemize}
\raggedright
  \item Say it once slowly, then at speed. $\to$ pronunciation reference
\end{itemize}
\end{sounds}

\begin{cousinweb}
\textbf{señora} — \textquotedblleft{}madam,\textquotedblright{} \textquotedblleft{}Mrs.\textquotedblright{}

The same \emph{senior}, turned by the \emph{-o}/\emph{-a} switch you have now met three times: \emph{bueno}/\emph{buena}, \emph{tranquilo}/\emph{tranquila}, and here.

That switch is the single most useful pattern in the language, and you have been collecting it since chapter one without being asked to memorise a rule.

One caution worth having early: \emph{señorita} exists and marks an unmarried woman. It is going out of use, and \emph{señora} is the safe address for an adult woman you do not know.
\end{cousinweb}

\begin{grammarlens}[title={What you've built}]
Both forms of address, and the \emph{-o}/\emph{-a} switch confirmed a third time.
\end{grammarlens}

\subsection*{Your turn}
Say these aloud:

\begin{itemize}
\raggedright
  \item \emph{señora}
  \item \emph{señora, por favor}
  \item \emph{señor — señora}
\end{itemize}

\begin{tcolorbox}[breakable,colback=teal!4,colframe=teal!35!black,title={Before you move on}]
What turns \emph{señor} into \emph{señora}? (The \emph{-o}/\emph{-a} switch.) Which is the safe address for an adult woman you do not know? (\emph{Señora}.)
\end{tcolorbox}
\section[encantado]{encantado — sung a spell over}

\label{lesson:ES-C305-encantado}



\begin{quote}
Two from last time, before anything new is added.
\end{quote}

\begin{sounds}
\begin{itemize}
\raggedright
  \item Say it once slowly, then at speed. $\to$ pronunciation reference
\end{itemize}
\end{sounds}

\begin{cousinweb}
\textbf{encantado} — what you say on being introduced.

Latin \textbf{incantare} is \emph{in-} plus \emph{cantare}, \textquotedblleft{}to sing\textquotedblright{} — to \textbf{sing a spell onto} something. English took it as \emph{enchanted}, and kept the raw noun as \emph{incantation}.

So the polite thing a Spanish speaker says on meeting you is, literally, \textquotedblleft{}I have been put under a spell.\textquotedblright{} English does the identical thing and nobody notices: \emph{charmed}, and \emph{delighted} — \emph{de-} plus \emph{lacere}, \textquotedblleft{}to lure.\textquotedblright{}

A man says \emph{encantado}; a woman says \emph{encantada}. You already have \emph{mucho gusto}, which does the same job more plainly — either is right.
\end{cousinweb}

\begin{grammarlens}[title={What you've built}]
A full introduction now: greet, give both names, address them properly, and be charmed.
\end{grammarlens}

\subsection*{Your turn}
Say these aloud:

\begin{itemize}
\raggedright
  \item \emph{encantado}
  \item \emph{encantada}
  \item \emph{me llamo… encantado}
\end{itemize}

\begin{tcolorbox}[breakable,colback=teal!4,colframe=teal!35!black,title={Before you move on}]
Take \emph{incantare} apart. (To sing a spell onto.) Which English word is it exactly? (\emph{Enchanted}.) Which word you already had does the same job? (\emph{Mucho gusto}.)
\end{tcolorbox}
